% frontmatter/preface.tex
\chapter{序言}

先看一个教学场景：一个Agent修好了栈的弹出函数，取出栈顶之后，还要从容器中移除它。任务结束，系统把这次经验留了下来。下一次遇到“只查看栈顶”的任务，它又读到了这条经验。此时，经验能否帮助行动，取决于条目有没有保留“需要改变容器内容”这个适用条件，也取决于检索器为什么选中它、执行器怎样理解它，以及最终检查能否发现误用。

这个例子连接了自进化系统的几个核心问题。经验要成为下一次行动的依据，需要经过提炼、保留、读取、执行和检验。每一步都有可以设计的机制，也有可以追踪的失败。沿着这条链研究，才能解释系统究竟学到了什么、为此花了多少资源，以及换一个任务后还能使用什么。

本书讨论\emph{经验如何变成可检验的能力变化}。读者将从状态与评价方法出发，进入记忆、提示、技能、工作流、参数学习和代码自修改，再研究更新时机与可靠性。机制解释和证据审计在这里共同构成一条主线：前者说明变化怎样发生，后者确定变化能够支持多大的结论。

\section*{用六问连接机制}

面对一个持续更新的系统，可以先把它还原成六个具体问题。
\begin{enumerate}[label=Q\arabic*.]
  \item \textbf{改什么？}指出提示、经验、技能、工作流或参数中实际改变的工件。
  \item \textbf{根据什么改？}追到任务轨迹、执行反馈、评审或奖励的来源。
  \item \textbf{怎样改？}说明候选如何生成、筛选并进入下一版本。
  \item \textbf{什么时候改？}区分任务内动作、任务间更新与离线训练。
  \item \textbf{保留什么？}列出跨任务继承的内容、读取规则及撤销边界。
  \item \textbf{怎样检验？}指定比较对象、共同资源、未参与选择的评价与失败判据。
\end{enumerate}

六问把方法名称变成可以检查的对象。以栈经验为例，“改什么”是新增条目，“保留什么”还包括条目的适用条件与读取规则；“怎样检验”则要同时查看弹出与只读两类任务。若随后把经验编成可执行技能，生成、验证、调用和维护的费用也随之改变。全书使用相同问题比较不同路线，让读者看见它们的共同结构和实际差别。

仍以栈为例。一次读取和一次改写各是求解过程中的一步，修复弹出函数的整次尝试构成一个回合。即使这次通过了测试，下一回合仍可能从原来的系统版本出发。只有当提炼出的经验被采纳、能被后续任务读取时，系统才带着这项变化进入下一版本。离线训练也可以改变后续任务的起点，此时版本中换入的是训练后的参数。第1、2章用步、回合和版本描述这几件事，第3章再问：保留这次变化，给后续任务带来了多少收益与成本？

\section*{沿证据理解收益与代价}

一个有用的结果应能从结论回到比较，再回到实际产生结果的工件。经验来自哪个任务，生成时看到了哪些信息，后续回合实际读取了什么，最终判定怎样形成，这些关系决定一个差值的含义。多次尝试、候选选择与持久更新都可能提高成绩，清楚记录它们，才能比较资源应当怎样分配。

成本也有时间结构。发现一个好策略可能昂贵，部署后却能反复使用；一次更新费用低，也可能增加每个后续任务的读取负担。书中的成本推导因此连接构建费、使用费、维护费与质量约束。读者可以据此计算在什么条件下值得采用一个方案，并指出哪些输入仍需测量。

全书以机制解释、案例审计和研究设计连接理论与实践。公开论文呈现不同方法在具体条件下的证据；编程试验提供可追溯的行为观察；教学计算展开可以复算的比较口径。第15章示范怎样沿归档理解结果，第16章把这些认识组织为十个研究问题，帮助读者选择下一项有价值的对照。问题沿用P编号，原研究假设与历史记录保存在附录\ref{sec:appg-history}；E03记录差异作为第5章的审计案例集中说明。

\section*{读完之后可以做什么}

本书面向研究生、算法工程师与Agent系统开发者，假定读者了解Python、语言模型调用和基础概率统计。附录C补充强化学习基础；附录B、D给出平台与复核入口，附录F展开成本账，附录G提供逐题导航和参考分析。

研究读者可以从估计量、对照和迁移范围进入，工程读者可以从工件、资源、准入和恢复进入。两条路径最终汇到一次具体的更新决策：采用什么改动，依据哪些结果，部署时保留哪些条件，出现什么观察就重新评价。书中提供的概念、推导与案例，目的都是帮助读者把这份决策写清楚，并沿实际记录检验它。
